\begin{Lemma}
\label{lem:adjoint-auxiliary-form-global-solvability}
The equation
\begin{equation}\label{eq:adjoint-auxiliary-form-global-solvability-equation}
[u,v]=\int_\Omega fv\,\mathrm{d}\bm{x}=(f,v)\quad\forall v\in H^{2m}(\Omega)
\end{equation}
has a solution $u\in H^{2m}(\Omega)$ if and only if $f\in(N^*)^\perp$. $u$ is unique up to addition of an arbitrary element of $N^*$.
\end{Lemma}

The following is a regularity theorem along the lines of the Agmon--Douglis--Nirenberg result.

\begin{Theorem}
\label{thm:adjoint-auxiliary-solution-higher-regularity}
Let $u$ be a solution of \eqref{eq:adjoint-auxiliary-form-global-solvability-equation}. Then $u\in H^{4m}(\Omega)$.
\end{Theorem}

The proof is hard and tedious and will not be given. We are now ready to state the main result of this subsection.

\begin{Theorem}
\label{thm:elliptic-adjoint-orthogonality-solvability}
The boundary-value problem
\begin{equation}\label{eq:elliptic-homogeneous-boundary-fredholm-solvability}
Lu=f,\ B_ju=0\text{ on }\partial\Omega\quad\text{for }j=1,\ldots,m,
\end{equation}
with $f\in L^2(\Omega)$ has a solution $u\in H^{2m}(\Omega)$ if and only if $f\in(N^*)^\perp$.
\end{Theorem}

\begin{Proof}
It is obvious from Green\textquotesingle s formula that the condition $f\in(N^*)^\perp$ is necessary. Let now $f\in(N^*)^\perp$. Then we can find $g\in M^*\cap H^{4m}(\Omega)$ such that $[g,v]=(f,v)$ for every $v\in H^{2m}(\Omega)$. Now let $u=L^*g\in H^{2m}(\Omega)$. For every $v\in H^{2m}(\Omega)$ which satisfies the boundary conditions $C_jv=0$, we conclude
\begin{equation}\label{eq:adjoint-solvability-proof-weak-interior-identity}
\int_\Omega u(L^*v)\,\mathrm{d}\bm{x}=\int_\Omega fv\,\mathrm{d}\bm{x}
\end{equation}
and, after integration by parts,
\begin{equation}\label{eq:adjoint-solvability-proof-recovered-boundary-identity}
\int_\Omega(Lu)v\,\mathrm{d}\bm{x}+\sum_{j=1}^m\int_{\partial\Omega}B_juT_jv\,\mathrm{d}S=\int_\Omega fv\,\mathrm{d}\bm{x}.
\end{equation}
It follows readily that $u$ satisfies $Lu=f$ and $B_ju=0$.
\end{Proof}

\begin{Remark}
\label{rem:adjoint-boundary-nullspace-uniqueness}
The adjoint boundary operators $C_j$ are generally not unique. However, it is clear from the last theorem that the space $N^*$ is uniquely determined. Hence different sets of adjoint boundary conditions are equivalent in the sense that they determine the same nullspace.
\end{Remark}
